% 2 câu đã gỡ khỏi ngân hàng (chuyen_dang) — chép lại vào file dạng để khôi phục

% từ: cau-hoi/lop12/2D1_Ung_dung_dao_ham_de_khao_sat_ham_so/2D1-3-1_GTLN_GTNN_tren_doan.tex
\begin{ex}%[2D1H3-1]%[ma:726bf6ec]
	Cho hàm số $y =  \dfrac{3(x^2+1)(x-3)}{(x+1)^3}$.
	 \choiceTF
		{Đồ thị hàm số không  có tiệm cận ngang}
		{\True Đồ thị hàm  số có đúng một tiệm cận ngang}
		{Đồ thị hàm số có một tiệm cận xiên} 
		{Đồ  thị hàm số không có tiệm cận đứng}
	 \loigiai{
		\begin{itemchoice}
			\itemch Ta có $\lim\limits_{x \to \pm \infty} \dfrac{3(x^2+1)(x-3)}{(x+1)^3} = \lim\limits_{x \to \pm \infty} \dfrac{3\left(1+\dfrac{1}{x^2} \right) \left(1-\dfrac{3}{x} \right)}{1 + \dfrac{3}{x}  + \dfrac{3}{x^2} + \dfrac{1}{x^3}} = 3$.\\
			Đồ thị hàm số có đúng một tiệm cận  ngang $y=3$.
			\itemch Đồ thị hàm số có đúng một tiệm cận  ngang.
			\itemch Đồ thị hàm số có tiệm cận ngang  thì không có tiệm cận xiên.
			\itemch Ta có  $\heva{&\lim\limits_{x \to -1^-}\dfrac{3(x^2+1)(x-3)}{(x+1)^3} = +\infty\\&
			\lim\limits{x\to -1^+}\dfrac{3(x^2+1)(x-3)}{(x+1)^3}  = +\infty.}$\\
			Vậy Đồ thị hàm số có một tiệm cận đứng  $x=-1$.
		\end{itemchoice}
	} 
\end{ex}

% từ: cau-hoi/lop12/2D3_Cac_so_dac_trung_ghep_nhom/2D3-2-1_Cong_thuc_ly_thuyet.tex
\begin{bt}%[2D3B2-1]%[ma:a7fcebe9]
	Cho $\displaystyle \int_{-1}^{3} f(x)\mathrm{\,d}x  = 7$. Tính $\displaystyle \int_{-1}^{3} [f(x)+4x]\mathrm{\,d}x$.
	 \shortans{23}
	\loigiai{ 
		Ta có $\displaystyle \int_{-1}^{3} [f(x)+4x]\mathrm{\,d}x = \displaystyle \int_{-1}^{3} [f(x)]\mathrm{\,d}x + \displaystyle  4\int_{-1}^{3} x\mathrm{\,d}x = 7 + 16 =23 $.
	} 
\end{bt}
